\section{Method}\label{sec:method}

\ifdefined\lemma\else\newtheorem{lemma}{Lemma}\fi

\subsection{Task and readout formats}\label{sec:method-task}

We use the belief-tracking stories of \citet{anonymous2026fitted}, generated by a port of their released template (the two ``initially sees'' sentences appear in alphabetical order of object name):
``Everyone initially sees that the \{$o$\} is in the \{init\}. Everyone initially sees that the \{$d$\} is in the \{dloc\}. \{$A$\} watches as the \{$o$\} is moved to the \{\textsc{loc}\}. \{$B$\} does not see this happen.'',
followed, in all main experiments, by the direct question ``Where does \{$A$\} believe the \{$o$\} is?''.
Locations are drawn from six single-token words (box, basket, shelf, drawer, cabinet, closet).
A base story B and a source story S differ only at the \emph{critical token} \{\textsc{loc}\}, at position $p$; we call it the state token.
In the cores we generate, the base, source and initial locations are distinct, and the distractor's location differs from the base and source locations.
For the identity test we use a third location X that is none of the base, source, initial or distractor locations, so the base prompt names it only in the candidate listing or the \fmtPost{} sentence.
We keep a core only if B, S and X tokenise identically except at $p$.

\paragraph{Encoder.} Except for one exploratory attention probe (\cref{app:additional}), every prompt is wrapped as in the \texttt{Answer:}-prefilled interface of \citet{anonymous2026fitted}: a system turn ``You are a helpful assistant.'', a user turn holding the raw prompt, the generation prompt with thinking disabled, and an ``Answer:'' assistant prefill.
We read full-vocabulary log-probabilities at the answer position and score the six candidate tokens (letters under \fmtLetter{}).

\paragraph{Format arms.} The story and question are identical across arms; only the candidate listing and the instruction change (full prompts in \cref{app:prompts}).
\begin{itemize}\setlength\itemsep{1pt}
  \item \fmtOpt{} (the format of \citealt{anonymous2026fitted}): after the question, ``Choices: box, basket, shelf, drawer, cabinet, closet'' and ``Answer with exactly one choice.''
  \item \fmtLetter{}: after the question, ``Choices: A) box, B) basket, \dots, F) closet'' and ``Answer with one letter.''; the answer is a letter.
  \item \fmtPost{}: the sentence ``The room has a box, a basket, a shelf, a drawer, a cabinet and a closet.'' is appended to the story; ``Answer with one word.''
  \item \fmtNone{}: no candidates; ``Answer with one word.''
  \item \fmtBefore{}: the choices line precedes the story; ``Answer with one word.'' Under the causal mask, the listed words cannot attend to the state token.
\end{itemize}
A \emph{re-mention} is any token after $p$ that names a candidate. \fmtOpt{}, \fmtLetter{} and \fmtPost{} re-mention all six locations after $p$; \fmtNone{} and \fmtBefore{} do not.

\subsection{Clamping the key and value channels}\label{sec:method-clamp}

Consider a decoder with causal self-attention. At layer $l$, position $t$ forms the query $q_t^l$ and reads earlier positions $u \le t$ only through their keys $K_u^l$ and values $V_u^l$. Hence a position $t > p$ can be influenced by the state token only through $\{K_p^l, V_p^l\}_l$: the key decides how much $t$ attends to $p$, and the value decides what is copied when it does. We call these the \emph{key channel} and the \emph{value channel}.

For a source run S, the clamp $\mathcal{C}_K(S)$ runs the base prompt but replaces, at position $p$ and every layer $l \ge l_0$, the output of \texttt{k\_proj} with the one cached from the S run; $\mathcal{C}_V(S)$ does the same for \texttt{v\_proj}, and $\mathcal{C}_{KV}(S)$ does both. The clamp acts on the \texttt{k\_proj} output, before any key normalisation (Qwen3, OLMo-2) and before RoPE; since both runs apply the same per-position normalisation and rotation at $p$, this equals clamping the final key. Under grouped-query attention, all KV groups are clamped. Our primary choice is $l_0 = 0$. We also use a depth-matched $l_0 = \mathrm{round}(0.0625\,L)$ (0-based) for $L$ layers; at Qwen2.5-72B ($L = 80$) this is 1-based block 6, the first exchanged block of \citet{anonymous2026fitted}.

\begin{lemma}[Exactness]\label{lem:exact}
If B and S differ only at position $p$, then $\mathcal{C}_{KV}(S)$ with $l_0 = 0$ reproduces the S run's hidden states at every position $t > p$, and hence its output distribution.
\end{lemma}
\emph{Proof.} Positions $u < p$ have identical inputs and, by causality, identical keys and values. Position $p$'s keys and values equal S's at every layer by construction. By induction over layers, every $t > p$ receives identical inputs and identical keys and values at all $u \le t$. \hfill$\square$

The state token's own residual stream is not reproduced, but no later position reads it except through the clamped channels. We verify the lemma numerically (\cref{app:impl}). It implies that the key and value channels together carry the entire effect of the state token on the answer.

\paragraph{Effects.} Let $m = \log p(\mathrm{S}) - \log p(\mathrm{B})$ at the answer position. For a clamp $C$, $d_C = m(C) - m(\mathcal{C}_{KV}(\mathrm{B}))$, where the identity row $\mathcal{C}_{KV}(\mathrm{B})$ clamps the base run to its own keys and values in the same batch, so that batch-dependent numerics cancel. We report $d_K$, $d_V$ and $d_{KV}$ for $\mathcal{C}_K(S)$, $\mathcal{C}_V(S)$ and $\mathcal{C}_{KV}(S)$. The \emph{key share} is $s_K = \bar d_K / (\bar d_K + \bar d_V)$, a ratio of means over cores, and the \emph{interaction} is $d_{KV} - d_K - d_V$.

\paragraph{Identity carried by keys.} A foreign key at $p$ can move $m$ without carrying the identity of S, for example by weakening the base answer, which raises all other candidates. To isolate identity we clamp to S or to X and measure absolute log-probability changes $\Delta\!\log p_Y(C) = \log p(Y \mid C) - \log p(Y \mid \mathcal{C}_{KV}(\mathrm{B}))$:
\begin{equation}\label{eq:idk}
\begin{split}
\mathrm{ID}_K = \tfrac{1}{2}\Big[ &\big(\Delta\!\log p_{\mathrm{S}}(K_{\mathrm{S}}) - \Delta\!\log p_{\mathrm{S}}(K_{\mathrm{X}})\big) \\
 + &\big(\Delta\!\log p_{\mathrm{X}}(K_{\mathrm{X}}) - \Delta\!\log p_{\mathrm{X}}(K_{\mathrm{S}})\big)\Big],
\end{split}
\end{equation}
where $K_Z$ denotes $\mathcal{C}_K(Z)$ with $l_0 = 0$. $\mathrm{ID}_V$ is defined analogously with values. If the effect of a foreign key decomposes as $\Delta\!\log p_Y(K_Z) = a_Y + b_Y\,\mathbb{1}[Y = Z]$, then $\mathrm{ID}_K = (b_{\mathrm{S}} + b_{\mathrm{X}})/2$. Any effect $a_Y$ shared by all foreign keys, such as weakening B or disturbing attention at $p$, cancels, and only the gain specific to the location the key came from remains. $\mathrm{ID}_K$ is measured in nats and is zero if keys carry no location identity.

\subsection{Localising the reader}\label{sec:method-splice}

To ask \emph{which} later tokens read the key, we make an exact row-restricted splice. At every layer the attention block runs twice on the same hidden states: once with the base key at $p$ and once with the source key $K_{\mathrm{S}}^l$ there. Output row $t$ takes the source version if $t \in G$ and the base version otherwise. Only the query rows in group $G$ thus see the swapped key; all other computation, including later reads of $G$'s hidden states, proceeds normally.
A unit test (Qwen2.5-0.5B, FP32) confirms that $G = $ all rows reproduces the full key clamp $\mathcal{C}_K(S)$ and $G = \emptyset$ the clean base run, each to within $10^{-4}$ in the logits; every run also re-checks $G = \emptyset$ (\cref{app:impl}).
The groups are the state token itself, the story tail after $p$, the question, the whole options line, the candidate words within it (or within the \fmtPost{} sentence), and the remaining rows after $p$ outside all of these (instruction, chat tokens, prefill and answer position). We report $[m(G) - m(\mathrm{B})]/[m(\text{all}) - m(\mathrm{B})]$, the fraction of the full key effect recovered. Restricting the splice to candidate-word rows, we further limit the swap to windows of four layers or to a single KV group's slice of the key.
We run the splice at Qwen2.5-1.5B (FP32, $n = 40$) under \fmtOpt{} and \fmtLetter{}. \todo{Stage 3: splice at Qwen2.5-7B, Qwen2.5-14B and Mistral-7B ($n = 60$; \fmtOpt{}, \fmtPost{}, \fmtLetter{}).}

\subsection{Revisiting the interventions of \citet{anonymous2026fitted}}\label{sec:method-paper1}

\paragraph{Bases and patch.} We use the released rank-16 bases of \citet{anonymous2026fitted} for Mistral-Small-24B and Qwen2.5-72B: the learned pair-swap remap M, the unfitted PCA basis P and the intended source-transfer basis $f^*$, each with seeds 101--103. The target is $\mathrm{T} = \pi(\mathrm{S})$, where $\pi$ swaps the six locations in pairs. For a basis $U \in \mathbb{R}^{16 \times d}$ with orthonormal rows, each event-span position is patched at the output of (1-based) block 4 as
\begin{equation}\label{eq:patch}
h \leftarrow h_{\mathrm{B}} + U^{\top} U\,(h_{\mathrm{S}} - h_{\mathrm{B}}).
\end{equation}
The event span covers both event sentences. Following their engine, we compute in BF16 with transformers 5.9.0 and pinned model revisions.

\paragraph{Reproduction check.} As a preregistered gate, the argmax token of M and of P under \fmtOpt{} must agree with the saved native outputs of \citet{anonymous2026fitted} on at least 0.95 of (core, fit) items, per model (\cref{app:prereg}).

\paragraph{Key-only exchange.} As in their native design, the critical-token key (pre-RoPE \texttt{k\_proj} output, all KV groups) of every block from 1-based block 6 onward is replaced by the other patched run's key. Queries, values and residual states evolve freely, and attention outputs are not restored. \emph{Addition} gives the PCA run the learned keys ($\mathrm{P}{+}K_{\mathrm{M}}$); \emph{removal} gives the learned run the PCA keys ($\mathrm{M}{+}K_{\mathrm{P}}$). Unlike our clamp, this exchange is not exact, because the patch also changes the other event-span tokens.

\paragraph{Metrics.} With $m = \log p(\mathrm{T}) - \log p(\mathrm{S})$,
\begin{align}
\varphi &= \frac{m(\mathrm{M}) - m(\mathrm{P})}{m(\mathrm{T}) - m(\mathrm{S})}, \label{eq:phi}\\
\psi &= \frac{m(\mathrm{P}{+}K_{\mathrm{M}}) - m(\mathrm{P})}{m(\mathrm{M}) - m(\mathrm{P})}, \label{eq:psi}\\
\rho &= \frac{m(\mathrm{M}{+}K_{\mathrm{P}}) - m(\mathrm{M})}{m(\mathrm{P}) - m(\mathrm{M})}. \label{eq:rho}
\end{align}
$\varphi$ is the learned remap's shift towards T relative to the natural S$\to$T shift, measured against the PCA patch. $\psi$ and $\rho$ are the fractions of the learned--PCA difference that the key-only exchange adds or removes. Each is computed as a ratio of means over cores, after averaging the three fits within each core.
For the preregistered transfer law we also use the value-channel completeness $\kappa_V$ of the learned remap. $\kappa_V$ was estimated before the runs from the fixed-value clamping data of \citet{anonymous2026fitted}, assuming that key and value effects add.

\paragraph{Population.} We use the 120 native story cores of \citet{anonymous2026fitted}. The reproduction gate uses all 120 (360 core--fit items per model); $\varphi$, $\psi$ and $\rho$ use the 96 cores in which B, S and T are distinct. In each of the other 24, B coincides with T or with S. Every core is evaluated under all five format arms.

\subsection{Models, statistics and preregistration}\label{sec:method-stats}

\paragraph{Models.} We study \nModels{} instruction-tuned models from four families: Qwen2.5-Instruct \citep{yang2024qwen25} at 1.5B, 3B, 7B, 14B, 32B and 72B; Qwen3-8B; Mistral-Small-24B-Instruct-2501 \citep{mistral2025small} and Mistral-7B-Instruct-v0.3; and OLMo-2-7B-Instruct \citep{olmo2025two}. GPU runs use BF16. Early tests at 1.5B ran on CPU in FP32.

\paragraph{Sample sizes and intervals.} The key/value factorial uses $n = 150$ story cores per format and model, with the direct question. The early CPU format factorials, preregistration A and the row-restricted splices at 1.5B used $n = 40$, and the GPU run repeated the factorial at 1.5B with $n = 150$. All intervals are 95\% percentile intervals from 10{,}000 bootstrap resamples of story cores. Ratios ($s_K$, $\varphi$, $\psi$, $\rho$) are recomputed as ratios of means within each resample. Splice fractions are ratios of means over cores and are reported as point estimates, without intervals. Items are pooled whether or not the model answers the clean runs correctly.

\paragraph{Preregistration.} We made four preregistrations (A--D; \cref{app:prereg}).
\begin{itemize}\setlength\itemsep{1pt}
  \item A: a fresh-seed CPU test at 1.5B of a claim formed after seeing earlier data.
  \item B: directional predictions in five out-of-sample open models at 3--8B (Qwen2.5-3B and 7B, Qwen3-8B, Mistral-7B, OLMo-2-7B), plus an exploratory, non-directional Qwen2.5 key-share curve from 1.5B to 14B.
  \item C: the released bases of \citet{anonymous2026fitted} at 24B and 72B, and the natural factorial at 24--72B.
  \item D: new tasks and localisation at 7--14B.
\end{itemize}
For each, the predictions were committed to the repository before the corresponding run, and the scoring script was committed before its outputs were inspected. We report every prediction of A--C as met or not met, and we label analyses outside the preregistrations as exploratory. D is pending. \todo{Stage 3: report the D1--D3 verdicts.} Before the GPU runs, an independent multi-agent review of the measurement code found several issues, which we fixed (\cref{app:impl}).
